\section{AIME II 2012}
        \begin{enumerate}
        	\item (\textit{AIME II 2012}) Find the number of ordered pairs of positive integer solutions $(m, n)$ to the equation $20m + 12n = 2012$.


 
 \textbf{Answer: }034



\textbf{Solution 1}

Solving for $m$ gives us $m = \frac{503-3n}{5},$ so in order for $m$ to be an integer, we must have $3n \equiv 503 \mod 5 \longrightarrow n \equiv 1 \mod 5.$ The smallest possible value of $n$ is obviously $1,$ and the greatest is $\frac{503 - 5}{3} = 166,$ so the total number of solutions is $\frac{166-1}{5}+1 = \boxed{034}$

 



\textbf{Solution 2}

Dividing by $4$ gives us $5m + 3n = 503$. 

 Thus, we have $3n = 503 -5m \implies 503 - 5m \equiv 0 \  (\text{mod }3)$, since $n$ is an integer.

 Rearranging:

 
 \[\begin{lgathered} 5m \equiv 503 \equiv 2 \  (\text{mod }3) \\ m \equiv 4 \  (\text{mod }3) \qquad \because 5^{-1} \equiv 2 \  (\text{mod }3)\\ m \equiv 1 \  (\text{mod }3) \\ \implies m = 3k + 1 \end{lgathered}\]

 Since $503 - 5m > 0$, we know that $m <\frac{503}{5}$.

 And because $m$ is an integer, $m \leq \left\lfloor  \frac{503}{5}  \right\rfloor \implies m \leq 100$

 Therefore, 

 
 \[\begin{lgathered} 3k + 1 \leq 100 \\ 0 \leq k \leq 33 \end{lgathered}\]

 Therefore, $m$ ranges from $3 \cdot 0+1$  to $3\cdot 33+1$, giving total $\boxed{ 034 }$ values.

 Alternatively, total integer values of $k$ is also $34$, thus, $\boxed{ 034 }$ values.

 Elaborated and reworded by Moisentinel

 



\textbf{Solution 3}

Because the x-intercept of the equation is $\frac{2012}{20}$, and the y-intercept is $\frac{2012}{12}$, the slope is $\frac{\frac{-2012}{12}}{\frac{2012}{20}} = \frac{-5}{3}$. Now, notice the first obvious solution: (100,1). From it, we derive all the other solutions by applying the slope in reverse, i.e: $(100,1), (97,6), (94,11)...$ Because the solutions are only positive, we can generate only 33 more solutions, so in total we have $33+1=\boxed{34}$ solutions.

 



\textbf{Solution 4 (Quick)}

Note that a positive integer is divisible by 20 if it ends in 0.

 Notice that if a multiple of 12 end in 2, 12 must be multiplied by an integer that ends in 1 or 6.

 So let's start checking because 2012 ends in 2, same as 12. 

 When $n=1$, $m=100$.

 When $n=6$, $m=97$.What is happening? Why doesn't, say, as values of $n$, do $1, 11, ...$ work, while $6, 16, 26$ work for $m$ to be an integer (as seen in 2022 CEMC Cayley \#21 (\href{https://www.cemc.uwaterloo.ca/contests/past\_contests/2022/2022CayleyContest.pdf}{https://www.cemc.uwaterloo.ca/contests/past\_contests/2022/2022CayleyContest.pdf}, \href{https://www.cemc.uwaterloo.ca/contests/past\_contests/2022/2022CayleySolution.pdf}{https://www.cemc.uwaterloo.ca/contests/past\_contests/2022/2022CayleySolution.pdf})?)

 Notice that we can break $1, 6, ...$ down into $0, 5, 10, ...$. So $20m=2000-12k$, where $k$ is a member of the set ${0, 5, 10, 15...}$.

 For a number to be divisible by 20, it must be divisible 10, and the quotient (that number divided by 10) must be divisible by 2. This means that the number must end in 0, and its tens digit must be even.

 So notice that the tens digit cycle in: $0, 4, 8, 2, 6$. Each of this is even (notice that when $k=25$, $2000-12k=1700$, it cycles again).

 So we know that all values of $n$, which end in 1 or 6, that make $2012-12n\geq20$ works.

 So values of $n$ that belong to ${1, 6, ..., 166}$ work. Clearly, that is 34 integer values of $n$. Therefore, the answer is $\boxed{034}$.

 


	\item (\textit{AIME II 2012}) Two geometric sequences $a_1, a_2, a_3, \ldots$ and $b_1, b_2, b_3, \ldots$ have the same common ratio, with $a_1 = 27$, $b_1=99$, and $a_{15}=b_{11}$. Find $a_9$.


 
 \textbf{Answer: }363



\textbf{Solution 1}

Call the common ratio $r.$ Now since the $n$th term of a geometric sequence with first term $x$ and common ratio $y$ is $xy^{n-1},$ we see that $a_1 \cdot r^{14} = b_1 \cdot r^{10} \implies r^4 = \frac{99}{27} = \frac{11}{3}.$ But $a_9$ equals $a_1 \cdot r^8 = a_1 \cdot (r^4)^2=27\cdot {\left(\frac{11}{3}\right)}^2=27\cdot \frac{121} 9=\boxed{363}$.

 



\textbf{Solution 2}

Let the ratio be \( r \). From \( \frac{a_{15}}{b_{11}} = \frac{a_{15}}{b_{11}} \):

 \( a_1 r^{14} = b_1 r^{10} \implies a_1 r^4 = b_1 \).

 Notice how \( a_5 = a_1 r^4 = b_1 \).

 Then

 \( a_9 = a_5 r^4 = b_1 r^4 = \frac{b_1^2}{a_1} \).

 Plug in \( a_1 = 27, b_1 = 99 \):

 \( a_9 = \frac{99^2}{27} = 363 \).

 $\boxed{363}$

 


	\item (\textit{AIME II 2012}) At a certain university, the division of mathematical sciences consists of the departments of mathematics, statistics, and computer science. There are two male and two female professors in each department. A committee of six professors is to contain three men and three women and must also contain two professors from each of the three departments. Find the number of possible committees that can be formed subject to these requirements.


 
 \textbf{Answer: }088



\textbf{Solution 1}

There are two cases:

 Case 1: One man and one woman is chosen from each department.

 Case 2: Two men are chosen from one department, two women are chosen from another department, and one man and one woman are chosen from the third department.

 For the first case, in each department there are ${{2}\choose{1}} \times {{2}\choose{1}} = 4$ ways to choose one man and one woman. Thus there are $4^3 = 64$ total possibilities conforming to case 1.

 For the second case, there is only ${{2}\choose{2}} = 1$ way to choose two professors of the same gender from a department, and again there are $4$ ways to choose one man and one woman. Thus there are $1 \cdot 1 \cdot 4 = 4$ ways to choose two men from one department, two women from another department, and one man and one woman from the third department. However, there are $3! = 6$ different department orders, so the total number of possibilities conforming to case 2 is $4 \cdot 6 = 24$.

 Summing these two values yields the final answer: $64 + 24 = \boxed{088}$.

 



\textbf{Solution 2}

Use \href{https://artofproblemsolving.com/wiki/index.php?title=Generating\_functions}{generating functions}. For each department, there is 1 way to pick 2 males, 4 ways to pick one of each, and 1 way to pick 2 females. Since there are three departments in total, and we wish for three males and three females, the answer will be equal to the coefficient of $x^3y^3$ in the expansion of $(x^2+4xy+y^2)^3$. The requested coefficient is $\boxed{088}$

 


	\item (\textit{AIME II 2012}) Ana, Bob, and Cao bike at constant rates of $8.6$ meters per second, $6.2$ meters per second, and $5$ meters per second, respectively. They all begin biking at the same time from the northeast corner of a rectangular field whose longer side runs due west. Ana starts biking along the edge of the field, initially heading west, Bob starts biking along the edge of the field, initially heading south, and Cao bikes in a straight line across the field to a point $D$ on the south edge of the field. Cao arrives at point $D$ at the same time that Ana and Bob arrive at $D$ for the first time. The ratio of the field's length to the field's width to the distance from point $D$ to the southeast corner of the field can be represented as $p : q : r$, where $p$, $q$, and $r$ are positive integers with $p$ and $q$ relatively prime. Find $p+q+r$.


 
 \textbf{Answer: }061



\textbf{Solution 1 (bash)}


                    \begin{figure}[H]
                        \centering
                        \includegraphics[width=0.5\textwidth]{Solution/AIME_II_2012_4_Solution_Figure_1.png}
                    \end{figure}
                    

 Let $a,b,c$ be the labeled lengths as shown in the diagram. Also, assume WLOG the time taken is $1$ second. 

 Observe that $\dfrac{2a+b+c}{8.6}=1$ or $2a+b+c=8.6$, and $\dfrac{b+c}{6.2}=1$ or $b+c=6.2$. Subtracting the second equation from the first gives $2a=2.4$, or $a=1.2$.

 Now, let us solve $b$ and $c$. Note that $\dfrac{\sqrt{b^2+c^2}}{5}=1$, or $b^2+c^2=25$. We also have $b+c=6.2$.

 We have a system of equations: 
 \[\left\{\begin{array}{l}b+c=6.2\\ b^2+c^2=25\end{array}\right.\]

 Squaring the first equation gives $b^2+2bc+c^2=38.44$, and subtracting the second from this gives $2bc=13.44$. Now subtracting this from $b^2+c^2=25$ gives $b^2-2bc+c^2=(b-c)^2=11.56$, or $b-c=3.4$. Now we have the following two equations:

 
 \[\left\{\begin{array}{l}b+c=6.2\\ b-c=3.4\end{array}\right.\]

 Adding the equations and dividing by two gives $b=4.8$, and it follows that $c=1.4$.

 The ratios we desire are therefore $1.4:6:4.8=7:30:24$, and our answer is $7+30+24=\boxed{061}$.

 Note that in our diagram, we labeled the part of the bottom $b$ and the side $c$. However, these labels are interchangeable. We can cancel out the case where the side is $4.8$ and the part of the bottom is $1.4$ by noting a restriction of the problem: "...a rectangular field whose longer side runs due west." If we had the side be $4.8$, then the entire bottom would be $1.2+1.4=2.6$, clearly less than $4.8$ and therefore violating our restriction.

 



\textbf{Solution 2 (uglier bash)}

Let P, Q, and R be the east-west distance of the field, the north-south distance, and the distance from the southeast corner to point D, respectively.

 Ana's distance to point D = $P + Q + (P - R) = 2P + Q - R$

 Bob's distance to point D = $Q + R$

 Cao's distance to point D = $\sqrt{Q^2 + R^2}$

 Since they arrive at the same time, their distance/speed ratios are equal, so:

 
 \[\frac{2P + Q - R}{8.6} = \frac{Q + R}{6.2} = \frac{\sqrt{Q^2 + R^2}}{5}\]

 
 \[\frac{2P + Q - R}{43} = \frac{Q + R}{31} = \frac{\sqrt{Q^2 + R^2}}{25}\]

 Looking at the last two parts of the equation:

 
 \[\frac{Q + R}{31} = \frac{\sqrt{Q^2 + R^2}}{25}\]

 
 \[25 (Q + R) = 31 \sqrt{Q^2 + R^2}\]

 
 \[625 Q^2 + 1250 QR + 625 R^2 = 961 Q^2 + 961 R^2\]

 
 \[336 Q^2 - 1250 QR + 336 R^2 = 0\]

 
 \[168 (\frac{Q}{R})^2 - 625 \frac{Q}{R}+ 168 = 0\]

 
 \[\frac{Q}{R} = \frac{625 \pm \sqrt{625^2 - 4 \cdot 168^2}}{336}\]

 
 \[\frac{Q}{R} = 24/7\;or\;7/24\]

 Looking at the first two parts of the equation above:

 
 \[\frac{2 P + Q - R}{43} = \frac{Q + R}{31}\]

 
 \[62 P + 31 Q - 31 R = 43 Q + 43 R\]

 
 \[P = \frac{6}{31}Q + \frac{37}{31} R\]

 If $\frac{Q}{R} = \frac{24}{7}$:

 
 \[R = \frac{7}{24} Q\]

 
 \[P = \frac{6}{31} Q + \frac{37}{31} \cdot \frac{7}{24} Q = \frac{13}{24} Q\]

 However, this makes P < Q, but we are given that P > Q.Therefore, $\frac{Q}{R} = \frac{7}{24}$, and:

 
 \[R = \frac{24}{7} Q\]

 
 \[P = \frac{6}{31} Q + \frac{37}{31} \cdot \frac{24}{7} Q = \frac{30}{7} Q\]

 
 \[P : Q : R = 30 : 7 : 24\]

 The solution is $P + Q + R = \framebox{061}$.

 


	\item (\textit{AIME II 2012}) In the accompanying figure, the outer square $S$ has side length $40$. A second square $S'$ of side length $15$ is constructed inside $S$ with the same center as $S$ and with sides parallel to those of $S$. From each midpoint of a side of $S$, segments are drawn to the two closest vertices of $S'$. The result is a four-pointed starlike figure inscribed in $S$. The star figure is cut out and then folded to form a pyramid with base $S'$. Find the volume of this pyramid.


 
                    \begin{figure}[H]
                        \centering
                        \includegraphics[width=0.5\textwidth]{Problem/AIME_II_2012_5_Problem_Figure_1.png}
                    \end{figure}
                    

 
 \textbf{Answer: }750



\textbf{Solution}

The volume of this pyramid can be found by the equation $V=\frac{1}{3}bh$, where $b$ is the base and $h$ is the height. The base is easy, since it is a square and has area $15^2=225$.

 
To find the height of the pyramid, the height of the four triangles is needed, which will be called $h^\prime$. By drawing a line through the middle of the larger square, we see that its length is equal to the length of the smaller rectangle and two of the triangle's heights. Then $40=2h^\prime +15$, which means that  $h^\prime=12.5$.

 
When the pyramid is made, you see that the height is the one of the legs of a right triangle, with the hypotenuse equal to $h^\prime$ and the other leg having length equal to half of the side length of the smaller square, or $7.5$. So, the Pythagorean Theorem can be used to find the height.

 $h=\sqrt{12.5^2-7.5^2}=\sqrt{100}=10$ 

 Finally, $V=\frac{1}{3}bh=\frac{1}{3}(225)(10)=\boxed{750}$.

 


	\item (\textit{AIME II 2012}) Let $z=a+bi$ be the complex number with $\vert z \vert = 5$ and $b > 0$ such that the distance between $(1+2i)z^3$ and $z^5$ is maximized, and let $z^4 = c+di$. Find $c+d$.


 
 \textbf{Answer: }125



\textbf{Solution}

Let's consider the maximization constraint first: we want to maximize the value of $|z^5 - (1+2i)z^3|$Simplifying, we have

 $|z^3| * |z^2 - (1+2i)|$

 $=|z|^3 * |z^2 - (1+2i)|$

 $=125|z^2 - (1+2i)|$

 Thus we only need to maximize the value of $|z^2 - (1+2i)|$.

 To maximize this value, we must have that $z^2$ is in the opposite direction of $1+2i$. The unit vector in the complex plane in the desired direction is $\frac{-1}{\sqrt{5}} + \frac{-2}{\sqrt{5}} i$. Furthermore, we know that the magnitude of $z^2$ is $25$, because the magnitude of $z$ is $5$. From this information, we can find that $z^2 = \sqrt{5} (-5 - 10i)$

 Squaring, we get $z^4 = 5 (25 - 100 + 100i) = -375 + 500i$. Finally, $c+d = -375 + 500 = \boxed{125}$

 



\textbf{Solution 2}

WLOG, let $z_{1}=5(\cos{\theta_{1}}+i\sin{\theta_{1}})$ and

 $z_{2}=1+2i=\sqrt{5}(\cos{\theta_{2}+i\sin{\theta_{2}}})$

 This means that

 $z_{1}^3=125(\cos{3\theta_{1}}+i\sin{3\theta_{1}})$

 $z_{1}^4=625(\cos{4\theta_{1}}+i\sin{4\theta_{1}})$

 Hence, this means that

 $z_{2}z_{1}^3=125\sqrt{5}(\cos({\theta_{2}+3\theta_{1}})+i\sin({\theta_{2}+3\theta_{1}}))$

 And

 $z_{1}^5=3125(\cos{5\theta_{1}}+i\sin{5\theta_{1}})$

 Now, common sense tells us that the distance between these two complex numbers is maxed when they both are points satisfying the equation of the line $yi=mx$, or when they are each a $180^{\circ}$ rotation away from each other.

 Hence, we must have that $5\theta_{1}=3\theta_{1}+\theta_{2}+180^{\circ}\implies\theta_{1}=\frac{\theta_{2}+180^{\circ}}{2}$

 Now, plug this back into $z_{1}^4$(if you want to know why, reread what we want in the problem!)

 So now, we have that$z_{1}^4=625(\cos{2\theta_{2}}+i\sin{2\theta_{2}})$

 Notice that $\cos\theta_{2}=\frac{1}{\sqrt{5}}$ and $\sin\theta_{2}=\frac{2}{\sqrt{5}}$

 Then, we have that $\cos{2\theta_{2}}=\cos^2{\theta_{2}}-\sin^2{\theta_{2}}=-\frac{3}{5}$ and $\sin{2\theta_{2}}=2\sin{\theta_{2}}\cos{\theta_{2}}=\frac{4}{5}$

 Finally, plugging back in, we find that $z_{1}^4=625(-\frac{3}{5}+\frac{4i}{5})=-375+500i$

 $-375+500=\boxed{125}$

 


 



\textbf{Solution 3}

Clearly, we want $\arg((1+2i)z^3) = \pi + \arg(z^5)$. This is equivalent to $\arg(1+2i) = \pi + 2 \arg(z)$ by the additive property of the argument in complex multiplication. Because we want to find $z^4$, we want an expression for $\arg(z^4)$. We now have $2\arg(1+2i) = 2\pi + 4 \arg(z) \rightarrow \arg(-3+4i) = 4\arg(z)$. Thus, $z^4$ is in the direction of $-3 + 4i$. To achieve a magnitude of 5\textasciicircum{}4, we need to multiple $-3+4i$ by 125, so $z^4 = 125(-3+4i)$ and the answer is thus $\boxed{125}$.

 


	\item (\textit{AIME II 2012}) Let $S$ be the increasing sequence of positive integers whose binary representation has exactly $8$ ones. Let $N$ be the 1000th number in $S$. Find the remainder when $N$ is divided by $1000$.


 
 \textbf{Answer: }032



\textbf{Solution}

Okay, an exercise in counting (lots of binomials to calculate!).  In base 2, the first number is $11111111$, which is the only way to choose 8 1's out of 8 spaces, or $\binom{8}{8}$.  What about 9 spaces?  Well, all told, there are $\binom{9}{8}=9$, which includes the first 1.  Similarly, for 10 spaces, there are $\binom{10}{8}=45,$ which includes the first 9.  For 11 spaces, there are $\binom{11}{8}=165$, which includes the first 45.  You're getting the handle.  For 12 spaces, there are $\binom{12}{8}=495$, which includes the first 165; for 13 spaces, there are $\binom{13}{8}=13 \cdot 99 > 1000$, so we now know that $N$ has exactly 13 spaces, so the $2^{12}$ digit is 1.  

 Now we just proceed with the other 12 spaces with 7 1's, and we're looking for the $1000-495=505th$ number.  Well, $\binom{11}{7}=330$, so we know that the $2^{11}$ digit also is 1, and we're left with finding the $505-330=175th$ number with 11 spaces and 6 1's.  Now $\binom{10}{6}=210,$ which is too big, but $\binom{9}{6}=84.$  Thus, the $2^9$ digit is 1, and we're now looking for the $175-84=91st$ number with 9 spaces and 5 1's.  Continuing the same process, $\binom{8}{5}=56$, so the $2^8$ digit is 1, and we're left to look for the $91-56=35th$ number with 8 spaces and 4 1's.  But here $\binom{7}{4}=35$, so N must be the last or largest 7-digit number with 4 1's.  Thus the last 8 digits of $N$ must be $01111000$, and to summarize, $N=1101101111000$ in base $2$.  Therefore, $N = 8+16+32+64+256+512+2048+4096 \equiv 32 \pmod{1000}$, and the answer is $\boxed{032}$.

 


 


	\item (\textit{AIME II 2012}) The complex numbers $z$ and $w$ satisfy the system 
 \[z + \frac{20i}w = 5+i\]

 \[w+\frac{12i}z = -4+10i\] Find the smallest possible value of $\vert zw\vert^2$.


 
 \textbf{Answer: }040



\textbf{Solution}

Multiplying the two equations together gives us 
 \[zw + 32i - \frac{240}{zw} = -30 + 46i\] and multiplying by $zw$ then gives us a quadratic in $zw$: 
 \[(zw)^2 + (30-14i)zw - 240 =0.\] Using the quadratic formula, we find the two possible values of $zw$ to be $7i-15 \pm \sqrt{(15-7i)^2 + 240}$ = $6+2i,$ $12i - 36.$ The smallest possible value of $\vert zw\vert^2$ is then obviously $6^2 + 2^2 = \boxed{040}$.

 \textbf{Note} A key thing to note here is that $|zw|^2=|z|^2\cdot|w|^2,$ which can be proved as follows:

 Proof: Using the values for $z$ and $w$ that we used above, we get:

 \begin{align*}|zw|^2&=|(ac-bd)+i(bc+ad)|^2\\&=(ac-bd)^2+(bc+ad)^2\\&=a^2c^2+b^2d^2+b^2c^2+a^2d^2-2abcd+2abc\\&=a^2c^2+b^2d^2+b^2c^2+a^2d^2\end{align*}
Also, $|z|^2=a^2+b^2$ and $|w|^2=c^2+d^2$. Therefore: 
 \[|z|^2\cdot|w|^2=(a^2+b^2)(c^2+d^2)=a^2c^2+a^2d^2+b^2c^2+b^2d^2\]and our proof is complete.

 Now, also note that we found $\sqrt{416-210i}$ by letting $416-210i=(a-bi)^2$ and solving for $a$ and $b$ by considering real and imaginary parts. Then, we substitute that into $a-bi$ which is the value of $\sqrt{416-210i}$ and continue from there.

 \textbf{mathboy282}

 


	\item (\textit{AIME II 2012}) Let $x$ and $y$ be real numbers such that $\frac{\sin x}{\sin y} = 3$ and $\frac{\cos x}{\cos y} = \frac12$. The value of $\frac{\sin 2x}{\sin 2y} + \frac{\cos 2x}{\cos 2y}$ can be expressed in the form $\frac pq$, where $p$ and $q$ are relatively prime positive integers. Find $p+q$.


 
 \textbf{Answer: }107



\textbf{Solution}

Examine the first term in the expression we want to evaluate, $\frac{\sin 2x}{\sin 2y}$, separately from the second term, $\frac{\cos 2x}{\cos 2y}$.

 \textbf{The First Term} Using the identity $\sin 2\theta = 2\sin\theta\cos\theta$, we have:

 $\frac{2\sin x \cos x}{2\sin y \cos y} = \frac{\sin x \cos x}{\sin y \cos y} = \frac{\sin x}{\sin y}\cdot\frac{\cos x}{\cos y}=3\cdot\frac{1}{2} = \frac{3}{2}$

 \textbf{The Second Term} Let the equation $\frac{\sin x}{\sin y} = 3$ be equation 1, and let the equation $\frac{\cos x}{\cos y} = \frac12$ be equation 2.Hungry for the widely-used identity $\sin^2\theta + \cos^2\theta = 1$, we cross multiply equation 1 by $\sin y$ and multiply equation 2 by $\cos y$.

 Equation 1 then becomes:

 $\sin x = 3\sin y$.

 Equation 2 then becomes:

 $\cos x = \frac{1}{2} \cos y$

 Aha! We can square both of the resulting equations and match up the resulting LHS with the resulting RHS:

 $1 = 9\sin^2 y + \frac{1}{4} \cos^2 y$

 Applying the identity $\cos^2 y = 1 - \sin^2 y$ (which is similar to $\sin^2\theta + \cos^2\theta = 1$ but a bit different), we can change $1 = 9\sin^2 y + \frac{1}{4} \cos^2 y$ into:

 $1 = 9\sin^2 y + \frac{1}{4} - \frac{1}{4} \sin^2 y$

 Rearranging, we get $\frac{3}{4} = \frac{35}{4} \sin^2 y$.

 So, $\sin^2 y = \frac{3}{35}$.

 Squaring Equation 1 (leading to $\sin^2 x = 9\sin^2 y$), we can solve for $\sin^2 x$:

 $\sin^2 x = 9\left(\frac{3}{35}\right) = \frac{27}{35}$

 Using the identity $\cos 2\theta = 1 - 2\sin^2\theta$, we can solve for $\frac{\cos 2x}{\cos 2y}$.

 $\cos 2x = 1 - 2\sin^2 x = 1 - 2\cdot\frac{27}{35} = 1 - \frac{54}{35} = -\frac{19}{35}$

 $\cos 2y = 1 - 2\sin^2 y = 1 - 2\cdot\frac{3}{35} = 1 - \frac{6}{35} = \frac{29}{35}$

 Thus, $\frac{\cos 2x}{\cos 2y} = \frac{-19/35}{29/35} = -\frac{19}{29}$.

 Plugging in the numbers we got back into the original equation :

 We get $\frac{\sin 2x}{\sin 2y} + \frac{\cos 2x}{\cos 2y} = \frac32 + \left(-\frac{19}{29} \right) = \frac{49}{58}$.

 So, the answer is $49+58=\boxed{107}$.

 



\textbf{Solution 2}

As mentioned above, the first term is clearly $\frac{3}{2}.$ For the second term, we first wish to find $\frac{\cos 2x}{\cos 2y} =\frac{2\cos^2 x - 1}{2 \cos^2y -1}.$ Now we first square the first equation getting $\frac{\sin^2x}{\sin^2y} =\frac{1-\cos^2x}{1 - \cos^2y} =9.$ Squaring the second equation yields $\frac{\cos^2x}{\cos^2y} =\frac{1}{4}.$ Let $\cos^2x = a$ and $\cos^2y = b.$ We have the system of equations
 \begin{align*} 1-a &= 9-9b \\ 4a &= b \\ \end{align*}
Multiplying the first equation by $4$ yields $4-4a = 36 - 36b$ and so $4-b =36 - 36b \implies b =\frac{32}{35}.$ We then find $a =\frac{8}{35}.$ Therefore the second fraction ends up being $\frac{\frac{64}{35}-1}{\frac{16}{35}-1} = -\frac{19}{29}$ so that means our desired sum is $\frac{49}{58}$ so the desired sum is $\boxed{107}.$

 



\textbf{Solution 3}


                    \begin{figure}[H]
                        \centering
                        \includegraphics[width=0.5\textwidth]{Solution/AIME_II_2012_9_Solution_Figure_1.png}
                    \end{figure}
                    

 We draw 2 right triangles with angles x and y that have the same hypotenuse.

 We get $b^2 + 9a^2 = 4b^2 + a^2$. Then, we find $8a^2 = 3b^2$.

 Now, we can scale the triangle such that $a = \sqrt{3}$, $b = \sqrt{8}$. We find all the side lengths, and we find the hypotenuse of both these triangles to equal $\sqrt{35}$ This allows us to find sin and cos easily.

 The first term is $\frac{3}{2}$, refer to solution 1 for how to find it.

 The second term is $\frac{\cos^2(x) - \sin^2(x)}{\cos^2(y) - \sin^2(y)}$. Using the diagram, we can easily compute this as $\frac{\frac{8}{35} - \frac{27}{35}}{\frac{32}{35} - \frac{3}{35}} = \frac{-19}{29}$

 Summing these you get $\frac{3}{2} + \frac{-19}{29} = \frac{49}{58} \implies \boxed{107}$

 -Alexlikemath

 



\textbf{Solution 4}

Let $a = \sin(x), b = \sin(y)$The first equation yields $\frac{a}{b} = 3.$ Using $sin^2(x) + cos^2(x) = 1$ the second equation yields 
 \[\frac{\sqrt{1-a^2}}{\sqrt{1-b^2}} = \frac{1}{2} \rightarrow \frac{1-a^2}{1-b^2} = \frac{1}{4}\]

 Solving this yields $\left(a, b\right) = \left(3\sqrt{\frac{3}{35}},\sqrt{\frac{3}{35}}\right).$Finding the first via double angle for sin yields
 \[\frac{\sin(2x)}{\sin(2y)} = \frac{2\sin{x}\cos{x}}{2\sin{y}\cos{y}} = 3 \cdot \frac{1}{2} = \frac{3}{2}\]Double angle for cosine is 
 \[\cos(2x) = 1-2\sin^2{x}\]so $\frac{\cos(2x)}{\sin(2x)} = \frac{1-2a^2}{1-2b^2} = -\frac{19}{29}.$Adding yields $\frac{49}{58} \rightarrow 49 + 58 = \boxed{107}$

 



\textbf{Solution 5}

We can calculate the first term $\frac{\sin 2x}{\sin 2y} = \frac{2 \sin x \cos x}{2 \sin y \cos y} = 3 \cdot \frac{1}{2} = \frac{3}{2}.$To calculate the second term, we need to use the identity $\sin^2 x + \cos^2 x = 1.$ From the first and second equations, we can rewrite then as $\sin x = 3 \sin y$ and $\cos x = \frac{1}{2} \cos y$ respectively. Now, we can use the identity and make the equation $9 \sin^2 y + \frac{1}{4} \cos^2 y = 1$ We now multiply both sides by 4 and get the equation $36 \sin^2 y + \cos^2 y = 4.$ Using the identity again, we realize that we can subtract 1 from both sides and obtain $35 \sin^2 y = 3$. Now, we can figure out that $\sin^2 y = \frac{3}{35}$. Another identity that is useful for this problem is $\cos 2x = 1 - 2 \sin^2 x$. Now, we can find the value of $\cos 2y$ 

 which is $1 - 2(\frac{3}{35}) = \frac{29}{35}$. Our main problem now is finding $\cos 2x$. We can use the identity so that we just need to find $\sin^2 x$. Using our first equation stated in the problem, we can multiply both sides by itself and get the equation $\frac{\sin^2 x}{\sin^2 y} = 9$. Plug in $\frac{3}{35}$ for $\sin^2 y$ and solve the equation to get the value of $\sin^2 x$ as $\frac{27}{35}$. Next, we use the identity and find $\cos 2x$ as $-\frac{19}{35}$. We can find the second term as we have $\cos 2x$ and $\cos 2y$. Thus, the total sum is $\frac{3}{2} + \frac{-19}{29} = \frac{49}{58}$. The question asks the sum of the numerator and the denominator so the answer is $49 + 58 = \boxed{107}$.

 



\textbf{Solution 6}

This is similar to solution 1.

 Define $a = \sin(x)$, $c = \sin(y)$, $b = \cos(x)$, $d= \cos(y)$. Therefore, $\frac{\sin(2x)}{\sin(2y)} = \frac{2ab}{2cd} = \frac{3}{2}$. 

 Now, finding $\frac{\cos(2x)}{\cos(2y)}$ will be harder. Firstly, through definition, $a = 3c$, $2b = d$. In addition, $a^2 + b^2 = 1$ and $c^2 + d^2 = 1$. So, $9c^2 + \frac{d^2}{4} = 1$ and $c^2 + d^2 = 1$. So, $36c^2 + d^2 = 4$ and $c^2 + d^2 = 1$. Therefore, $35c^2 = 3$, and so $c^2 = \frac{3}{35}$. Then, $d^2 = 1-\frac{3}{35} = \frac{32}{35}$. So, $4b^2 = \frac{32}{35}$ and $b^2 = \frac{8}{35}$. The identity $\cos(2x) = 2b^2 - 1$ holds, and similarly, $2d^2 - 1$ holds, so we have $\cos(2x) = -\frac{19}{35}$ and $\cos(2y) = \frac{29}{35}$. Thus the requested value $\frac{\cos(2x)}{\cos(2y)} = -\frac{19}{29}$.

 The requested value is $\frac{3}{2} - \frac{19}{29} = \frac{49}{58} \implies \boxed{107}$.

 


	\item (\textit{AIME II 2012}) Find the number of positive integers $n$ less than $1000$ for which there exists a positive real number $x$ such that $n=x\lfloor x \rfloor$.


 Note: $\lfloor x \rfloor$ is the greatest integer less than or equal to $x$.


 
 \textbf{Answer: }496



\textbf{Solution 1}

We know that $x$ cannot be irrational because the product of a rational number and an irrational number is irrational (but $n$ is an integer). Therefore $x$ is rational.

 
Let $x = a + \frac{b}{c}$ where $a,b,c$ are nonnegative integers and $0 \le b < c$ (essentially, $x$ is a mixed number). Then,
 \[n = \left(a + \frac{b}{c}\right) \left\lfloor a +\frac{b}{c} \right\rfloor \Rightarrow n = \left(a + \frac{b}{c}\right)a = a^2 + \frac{ab}{c}\] 

 Here it is sufficient for $\frac{ab}{c}$ to be an integer. We can use casework to find values of $n$ based on the value of $a$:

 $a = 0 \implies$ nothing because n is positive

 $a = 1 \implies \frac{b}{c} = \frac{0}{1}$

 $a = 2 \implies \frac{b}{c} = \frac{0}{2},\frac{1}{2}$

 $a = 3 \implies\frac{b}{c} =\frac{0}{3},\frac{1}{3},\frac{2}{3}$

 
The pattern continues up to $a = 31$. Note that if $a = 32$, then $n > 1000$. However if $a = 31$, the largest possible $x$ is $31 + \frac{30}{31}$, in which $n$ is still less than $1000$. Therefore the number of positive integers for $n$ is equal to $1+2+3+...+31 = \frac{31 \cdot 32}{2} = \boxed{496}.$

 



\textbf{Solution 2}

Notice that $x\lfloor x\rfloor$ is continuous over the region $x \in [k, k+1)$ for any integer $k$. Therefore, it takes all values in the range $[k\lfloor k\rfloor, (k+1)\lfloor k+1\rfloor) = [k^2, (k+1)k)$ over that interval. Note that if $k>32$ then $k^2 > 1000$ and if $k=31$, the maximum value attained is $31*32 < 1000$. It follows that the answer is $\sum_{k=1}^{31} (k+1)k-k^2  = \sum_{k=1}^{31} k = \frac{31\cdot 32}{2} = \boxed{496}.$

 



\textbf{Solution 3}

Bounding gives $x^2\le n<x^2+x$. Thus there are a total of $x$ possible values for $n$, for each value of $x^2$. Checking, we see $31^2+31=992<1000$, so there are 
 \[1+2+3+...+31= \boxed{496}\] such values for $n$.

 Here writer's x and question's x are different.
Here writer means [x] everywhere he writes x
so you can let [x] = k and then use the inequality $k^2\le n<k^2+k$.
this inequality comes from [x] = k
$k\le x<k+1$
$k[x]\le x[x]<(k+1)[x]$
$k^2\le n<(k+1)k$
 further you can proceed.

 



\textbf{Solution 4}

After a bit of experimenting, we let $n=l^2+s, s < 2n+1$. We claim that I (the integer part of $x$) = $l$ . (Prove it yourself using contradiction !) so now we get that $x=l+\frac{s}{l}$. This implies that solutions exist iff $s<l$, or for all natural numbers of the form $l^2+s$ where $s<l$. Hence, 1 solution exists for $l=1$! 2 for $l=2$ and so on. Therefore our final answer is $31+30+\dots+1= \boxed{496}$

 



\textbf{Solution 5}

Notice we can write $\lfloor x \rfloor$ as $x - \{x\}$. Since $n$ is a positive integer less than $1000$, 

 
 \[1 \leq x(x - \{x\}) \leq 999.\]

 We now have 

 
 \[1 \leq x^{2} - x\{x\} \leq 999.\]

 Since $x\{x\}$ ranges from $0$ to $x$, we can rewrite as 

 
 \[1 \leq x^{2} \leq 999 + x.\]

 Notice all $x$ that satisfy this inequality range from $1$ to $32$. 

 Now every range from $y$ to $y+1$ for $x$ results in exactly $y$ solutions that give a positive integer solution for $n$. The first range is $1$ to $2$ and the last range is $31$ to $32$. 

 Thus, by using our above lemma, we can find the total number of solutions:

 
 \[1+2+\cdots+31 = \frac{31 \cdot 32}{2} = \boxed{496}.\]

 


 



\textbf{Solution 6}

Let $x = a + \alpha$ where $a$ is a positive integer and $0 \leq \alpha < 1$. Then we have

 $(a + \alpha)(a) = n$ since $\lfloor a + \alpha \rfloor = a$ 

 Now we get

 $a^2 + a \alpha = n \implies \alpha = \frac{n - a^2}{a}$ so we get

 $0 \leq \frac{n - a^2}{a} < 1 \implies a^2 \leq n < a^2 + a$

 Since $n$ is an integer, it means all $n$ possible are between $a^2$ and $a^2 + a - 1$ inclusive. The number of possibilities here is $a^2 + a - 1 - a^2 + 1 = a$. Since $n < 1000$ we get $a^2 < 1000 \implies a \leq 31$. Therefore, for each $1 \leq a \leq 31$, as $a$ is a positive integer, we get $a$ solutions for $n$. In other words, we simply have

 $1 + 2 + \cdots + 31 = \boxed{496}$

 


	\item (\textit{AIME II 2012}) Let $f_1(x) = \frac23 - \frac3{3x+1}$, and for $n \ge 2$, define $f_n(x) = f_1(f_{n-1}(x))$. The value of $x$ that satisfies $f_{1001}(x) = x-3$ can be expressed in the form $\frac mn$, where $m$ and $n$ are relatively prime positive integers. Find $m+n$.


 
 \textbf{Answer: }008



\textbf{Solution}

After evaluating the first few values of $f_k (x)$, we obtain $f_4(x) = f_1(x) = \frac{2}{3} - \frac{3}{3x+1} = \frac{6x-7}{9x+3}$. Since $1001 \equiv 2 \mod 3$, $f_{1001}(x) = f_2(x) = \frac{3x+7}{6-9x}$. We set this equal to $x-3$, i.e.

 
$\frac{3x+7}{6-9x} = x-3 \Rightarrow x = \frac{5}{3}$. The answer is thus $5+3 = \boxed{008}$.

 



\textbf{Solution 2 (Solution 1, With work)}


 \[f_1(x) = \frac{2}{3} - \frac{3}{3x+1}\]
 \[f_2(x) = \frac{2}{3} - \frac{3}{3(\frac{2}{3} - \frac{3}{3x+1}) + 1}\]
 \[f_2(x) = \frac{2}{3} - \frac{3}{3 - \frac{9}{3x+1}}\]
 \[f_2(x) = \frac{2}{3} - 3 \cdot \frac{3x+1}{9x-6}\]
 \[f_2(x) = \frac{-3x-7}{9x-6}\]
 \[f_3(x) = \frac{2}{3} - \frac{3}{3\frac{-3x-7}{9x-6} + 1}\]
 \[f_3(x) = \frac{2}{3} - \frac{3}{\frac{-3x-7}{3x-2} + 1}\]
 \[f_3(x) = \frac{2}{3} - \frac{3}{\frac{-9}{3x-2}})\]
 \[f_3(x) = \frac{2}{3} - \frac{2-3x}{3}\]
 \[f_3(x) = \frac{3x}{3}\]
 \[f_3(x) = x\].

 Therefore, because $f_1(f_1(f_1(x))) = x$, it follows that $f_1(x)$ is cyclic with cyclicity 3.

 Because $1001 \equiv 2 \pmod{3}$, $f_{1001}(x) = x-3$, we have $x-3 = f_2(x) = \frac{-3x-7}{9x-6}$, so $(x-3)(9x-6) = -(3x+7)$. In other words, $9x^2 -30x+25 = 0$. Notice that $30 = 3 \cdot 2 \cdot 5$ so we factor into $(3x-5)^2 =0$, Then, $x = \frac{5}{3} \implies \boxed{8}$.

 


	\item (\textit{AIME II 2012}) For a positive integer $p$, define the positive integer $n$ to be $p$\textit{-safe} if $n$ differs in absolute value by more than $2$ from all multiples of $p$. For example, the set of $10$-safe numbers is $\{ 3, 4, 5, 6, 7, 13, 14, 15, 16, 17, 23, \ldots\}$. Find the number of positive integers less than or equal to $10,000$ which are simultaneously $7$-safe, $11$-safe, and $13$-safe.


 
 \textbf{Answer: }958



\textbf{Solution}

We see that a number $n$ is $p$-safe if and only if the residue of $n \mod p$ is greater than $2$ and less than $p-2$; thus, there are $p-5$ residues $\mod p$ that a $p$-safe number can have. Therefore, a number $n$ satisfying the conditions of the problem can have $2$ different residues $\mod 7$, $6$ different residues $\mod 11$, and $8$ different residues $\mod 13$. The Chinese Remainder Theorem states that for a number $x$ that is$a$ (mod b)$c$ (mod d)$e$ (mod f)has one solution if $\gcd(b,d,f)=1$. For example, in our case, the number $n$ can be:3 (mod 7)3 (mod 11)7 (mod 13)so since $\gcd(7,11,13)$=1, there is 1 solution for n for this case of residues of $n$.

 This means that by the Chinese Remainder Theorem, $n$ can have $2\cdot 6 \cdot 8 = 96$ different residues mod $7 \cdot 11 \cdot 13 = 1001$. Thus, there are $960$ values of $n$ satisfying the conditions in the range $0 < n \le 10010$. However, we must now remove any values greater than $10000$ that satisfy the conditions. By checking residues, we easily see that the only such values are $10006$ and $10007$, so there remain $\fbox{958}$ values satisfying the conditions of the problem.

 


 
                \begin{itemize}
                    \item We can also say $\frac{2}{7}$ of all numbers are 7-safe, $\frac{6}{11}$ of all numbers are 11-safe, and $\frac{8}{13}$ of all numbers are 13-safe. We can multiply these to get that $\frac{96}{1001}$ of all numbers are simultaneously 7-safe, 11-safe, and 13-safe.
                \end{itemize}


                



\textbf{Solution 2}

We want the number of integers $n\le 10000$ such that\begin{align*}n&\equiv 3,4\bmod{7}\\n&\equiv 3,4,5,6,7,8,\bmod{11}\\n&\equiv 3,4,5,6,7,8,9,10\bmod{13}\end{align*}
By the Chinese Remainder Theorem, the map
 \[f \colon \mathbb{Z}/1001\mathbb{Z} \to \mathbb{Z}/7\mathbb{Z} \times \mathbb{Z}/11\mathbb{Z} \times \mathbb{Z}/13\mathbb{Z},\]
 \[n \bmod 1001 \mapsto (n \bmod 7, n \bmod 11, n \bmod 13)\]is bijective. This means that the residues of $n$ modulo $7$, $11$, and $13$ uniquely determine its residue modulo $1001$. Hence, for $n\in [1,10010]$, there are $2\times 6\times 8\times 10=960$ integers that are simultaneously $7$-safe, $11$-safe, and $13$-safe.
                    \begin{figure}[H]
                        \centering
                        \includegraphics[width=0.5\textwidth]{Solution/AIME_II_2012_12_Solution_Figure_1.png}
                    \end{figure}
                    However, this accounts the case when $n=10006,10007$, which are out-of-range. Therefore, the total is $960-2=\boxed{958}$.

 


	\item (\textit{AIME II 2012}) Equilateral $\triangle ABC$ has side length $\sqrt{111}$. There are four distinct triangles $AD_1E_1$, $AD_1E_2$, $AD_2E_3$, and $AD_2E_4$, each congruent to $\triangle ABC$,
with $BD_1 = BD_2 = \sqrt{11}$. Find $\sum_{k=1}^4(CE_k)^2$.


 
 \textbf{Answer: }677



\textbf{Solution 1}

Let $\angle BAD_1 = \angle BAD_2 = \theta$. By law of cosines, 
 \[\sum_{k=1}^4(CE_k)^2 = 888-222\sum_{k=1}^4 \cos \angle CAE_k \quad (\ast)\] Consider the inside sum. Note that the triangles in question are just rotations of $\triangle ABC$ around $A$. This allows us to find 
 \[\angle CAE_1 = \theta, \angle CAE_2 = 120^\circ - \theta, \angle CAE_3 = \theta, \angle CAE_4 = 120^\circ + \theta\] Thus, applying sum to product, 
 \[\sum_{k=1}^4 \cos \angle CAE_k = \cos \theta +\cos \theta + \cos (120^\circ +\theta)+\cos (120^\circ-\theta) = 2\cos \theta +2\cos 120^\circ \cos \theta\] 
 \[= \cos \theta =\frac{(\sqrt{111})^2+(\sqrt{111})^2-(\sqrt{11})^2}{2 (\sqrt{111})(\sqrt{111})} = \frac{211}{222}\] Finally, $(\ast)$ becomes 
 \[\sum_{k=1}^4(CE_k)^2 = 888-222\sum_{k=1}^4 \cos \angle CAE_k = \boxed{677}\]

 -sillybone

 



\textbf{Solution 2}

Note that there are only two possible locations for points $D_1$ and $D_2$, as they are both $\sqrt{111}$ from point $A$ and $\sqrt{11}$ from point $B$, so they are the two points where a circle centered at $A$ with radius $\sqrt{111}$ and a circle centered at $B$ with radius $\sqrt{11}$ intersect.  Let $D_1$ be the point on the opposite side of $\overline{AB}$ from $C$, and $D_2$ the point on the same side of $\overline{AB}$ as $C$.

 Let $\theta$ be the measure of angle $BAD_1$ (also the measure of angle $BAD_2$); by the Law of Cosines,
 \begin{align*}\sqrt{11}^2 &= \sqrt{111}^2 + \sqrt{111}^2 - 2 \cdot \sqrt{111} \cdot \sqrt{111} \cdot \cos\theta\\ 11 &= 222(1 - \cos\theta)\end{align*}


 There are two equilateral triangles with $\overline{AD_1}$ as a side; let $E_1$ be the third vertex that is farthest from $C$, and $E_2$ be the third vertex that is nearest to $C$.

 Angle $E_1AC = E_1AD_1 + D_1AB + BAC = 60 + \theta + 60 = 120 + \theta$; by the Law of Cosines,
 \begin{align*}(E_1C)^2 &= (E_1A)^2 + (AC)^2 - 2 (E_1A) (AC)\cos(120 + \theta)\\ &= 111 + 111 - 222\cos(120 + \theta)\end{align*}
Angle $E_2AC = \theta$; by the Law of Cosines,
 \begin{align*}(E_2C)^2 &= (E_2A)^2 + (AC)^2 - 2 (E_2A) (AC)\cos\theta\\ &= 111 + 111 - 222\,\cos\theta\end{align*}


 There are two equilateral triangles with $\overline{AD_2}$ as a side; let $E_3$ be the third vertex that is farthest from $C$, and $E_4$ be the third vertex that is nearest to $C$.

 Angle $E_3AC = E_3AB + BAC = (60 - \theta) + 60 = 120 - \theta$; by the Law of Cosines,
 \begin{align*}(E_3C)^2 &= (E_3A)^2 + (AC)^2 - 2 (E_3A) (AC)\cos(120 - \theta)\\ &= 111 + 111 - 222\cos(120 - \theta)\end{align*}
Angle $E_4AC = \theta$; by the Law of Cosines,
 \begin{align*}(E_4C)^2 &= (E_4A)^2 + (AC)^2 - 2 (E_4A) (AC)\cos\theta \\ &= 111 + 111 - 222\cos\theta\end{align*}


 The solution is:
 \begin{align*} \sum_{k=1}^4(CE_k)^2 &= (E_1C)^2 + (E_3C)^2 + (E_2C)^2 + (E_4C)^2\\ &= 222(1 - \cos(120 + \theta)) + 222(1 - \cos(120 - \theta)) + 222(1 - \cos\theta) + 222(1 - \cos\theta)\\ &= 222((1 - (\cos120\cos\theta - \sin120\sin\theta)) + (1 - (\cos120\cos\theta + \sin120\sin\theta)) + 2(1 -\cos\theta))\\ &= 222(1 - \cos120\cos\theta + \sin120\sin\theta + 1 - \cos120\cos\theta - \sin120\sin\theta + 2 - 2\cos\theta)\\ &= 222(1 + \frac{1}{2}\cos\theta + 1 + \frac{1}{2}\cos\theta + 2 - 2\cos\theta)\\ &= 222(4 - \cos\theta)\\ &= 666 + 222(1 - \cos\theta) \end{align*}
Substituting $11$ for $222(1 - \cos\theta)$ gives the solution $666 + 11 = \framebox{677}.$

 



\textbf{Solution 3}

This problem is pretty much destroyed by complex plane geometry, which is similar to vector geometry only with the power of easy rotation.  Place the triangle in the complex plane by letting $C$ be the origin, placing $B$ along the x-axis, and $A$ in the first quadrant.  Let $r=\sqrt{111}$.  If $\omega$ denotes the primative sixth root of unity, $e^{i\pi/3}$, then we have $C=0$, $B=r$, and $A=r\omega.$  Recall that counter-clockwise rotation in the complex plane by an angle $\theta$ is accomplished by multiplication by $e^{i\theta}$ (and clockwise rotation is multiplication by its conjugate).  So, we can find $D_1$ and $D_2$ by rotating $B$ around $A$ by angles of $\theta$ and $-\theta$, where $\theta$ is the apex angle in the isoceles triangle with sides $\sqrt{111}$, $\sqrt{111}$, and $\sqrt{11}$.  That is, let $z=e^{i\theta}$, and then:

 $D_1=A+z(B-A)$, and $D_2=A+\overline{z}(B-A)$.  Now notice that $B-A=\overline{A}$, so this simplifies further to:

 $D_1=A+z\overline{A}$, and $D_2=A+\overline{z}\overline{A}$.  

 Similarly, we can write $E_1$, $E_2$, $E_3$, and $E_4$ by rotating $D_1$ and $D_2$ around $A$ by $\pm\pi/3$:

 $E_1=A+\omega(D_1-A)$, $E_2=A+\overline{\omega}(D_1-A)$, $E_3=A+\omega(D_2-A)$, $E_4=A+\overline{\omega}(D_2-A)$.  Thus:

 $E_1=A+\omega z \overline{A}$, $E_2=A+\overline{\omega} z \overline{A}$, $E_3=A+\omega\overline{z}\overline{A}$, $E_4=A+\overline{\omega}\overline{z}\overline{A}$.

 Now to find some magnitudes, which is easy since we chose $C$ as the origin:

 $\|E_1\|^2=(A+\omega z \overline{A})(\overline{A}+\overline{\omega z}A)=2\|A\|^2+\omega z \overline{A}^2 + \overline{\omega}\overline{z}A^2$,

 $\|E_2\|^2=(A+\overline{\omega} z \overline{A})(\overline{A}+\omega \overline{z}A)=2\|A\|^2+\overline{\omega} z \overline{A}^2 + \omega\overline{z}A^2$,

 $\|E_3\|^2=(A+\omega \overline{z} \overline{A})(\overline{A}+\overline{\omega} z A)=2\|A\|^2+\omega \overline{z} \overline{A}^2 + \overline{\omega}zA^2$,

 $\|E_4\|^2=(A+\overline{\omega z} \overline{A})(\overline{A}+\omega zA)=2\|A\|^2+\overline{\omega z} \overline{A}^2 + \omega zA^2$.

 Adding these up, the sum equals $8\|A\|^2+(\overline{A}^2+A^2)(\omega z + \omega \overline{z} + \overline{\omega}z + \overline{\omega}\overline{z}) = 8\|A\|^2+(\overline{A}^2+A^2)(z + \overline{z})( \omega+ \overline{\omega})$.  

 (Isn't that nice?)  Notice that $\overline{A}^2+A^2 = r^2(\overline{\omega}^2+\omega^2) = -r^2$, and $\omega+ \overline{\omega}=1$, so that this sum simplifies further to $888-111(z + \overline{z})$.  

 Finally, $z + \overline{z} = 2\cos{\theta}$, which is found using the law of cosines on that isoceles triangle: $11=111+111-222\cos{\theta}$, so $2\cos{\theta}=211/111$.

 Thus, the sum equals $888-211=\framebox{677}$.

 



\textbf{Solution 4}

This method uses complex numbers with $A$ as the origin. Let $A=0$, $B=\sqrt{111}$, $C = \sqrt{111}\theta$, where $\theta = e^{i \pi/3} = \frac{1}{2} + \frac{\sqrt{3}}{2}i$.

 Also, let $x$ be $D_1$ or $D_2$. Then

 $|x|=\sqrt{111}, |x-\sqrt{111}|=\sqrt{11}$

 Therefore, $11 = |x-\sqrt{111}|^2 = |x|^2 + 111 -2\sqrt{111}Re(x) = 222 - 2\sqrt{111}Re(x)$, so

 
 \[2\sqrt{111}Re(x) = 211.\]

 Since $E_1$, $E_2$ are one of $D_1\theta$ or $D_1\theta^{-1}$, without loss of generality, let $E_1=D_1\theta$ and $E_2 = D_1\theta^{-1}$. Then

 
 \[|CE_1|^2 = |\sqrt{111}-D_1|^2 = |\sqrt{111}-x|^2=11\]
 \[|CE_2|^2 = |\theta^2 \sqrt{111} - D_1|^2 = 222-2\sqrt{111} Re(D_1 \theta^2)\]

 One can similarly get $|CE_3|=11$ and $|CE_4|=222-2\sqrt{111} Re(D_2 \theta^2)$, so the desired sum is equal to

 
 \[22+444-2\sqrt{111}(Re(D_1 \theta^2)+Re(D_1 \theta^2))\]

 Note that $Re(D_{1,2} \theta^2) = Re((Re(x) \pm Im(x)i)(-1/2 + \sqrt{3}/2i)) = - (Re(x)/2 \pm Im(x)\sqrt{3}/2)$, so the sum of these two is just $-Re(x)$. Therefore the desired sum is equal to 

 
 \[22+444+2\sqrt{111}Re(x) = 22+444+211 = \framebox{677}.\]

 



\textbf{Solution 5}

This method uses the observation that every point is equidistant from $A$. Without loss of generality, we can assume $C$ is on the same side of $AB$ as $D_1$.

 We can start off by angle chasing the angles around $A$. We let $\angle BAD_1 = \alpha$. Then, we note that $BD_1 = BD_2$ and $AD_1 = AB = AD_2$. Thus, $\bigtriangleup ABD_1 \cong \bigtriangleup ABD_2$. Thus, $\angle BAD_2 = \alpha$ also.

 We can now angle chase the angles about $A$. Because $\angle D_2AE_3 = 60$, $\angle D_1AE_3 = 60- 2 \alpha$. We can use all the congruent equilateral triangles in a similar manner obtaining:
 \[\angle D_1AE_3 = 60- 2 \alpha\]
 \[\angle E_3AC = \alpha\]
 \[\angle CAE_1 = \alpha\]
 \[\angle D_2AE_2 = 60- 2 \alpha\]
 \[\angle E_2AE_4 = 2 \alpha\]

 Now, $AE_3 = AC = \sqrt{111}$ and $\angle E_3AC = \alpha$. Thus, $\bigtriangleup E_3AC \cong \bigtriangleup BAD_1$. Thus, $CE_3^2 = BD_1^2 = 11$.

 Similarly, $AC = AE_1 = \sqrt{111}$ and $\angle CAE_1 = \alpha$. Thus, $\bigtriangleup CAE_1 \cong \bigtriangleup BAD_1$. Thus, $CE_1^2 = BD_1^2 = 11$.

 We can use $\bigtriangleup CAE_2$ to find $CE_2^2$. Law of Cosines yields
 \[CE_2^2 = AE_2^2 + AC^2 - 2 \cdot AE_2 \cdot AC \cdot \cos(\angle E_2AC).\]Substituting the known lengths and angles gives
 \[CE_2^2 = 222 - 222 \cdot \cos(120- \alpha).\]Expanding this with the Cosine Subtraction Identity we get
 \[CE_2^2 = 222 - 222(\cos120 \cos \alpha + \sin120 \sin \alpha).\]

 We could attempt to calculate this but we can clear it up by simultaneously finding $CE_4^2$ too.We use Law of Cosines on $\bigtriangleup CAE_4$ to get
 \[CE_4^2 = AE_4^2 + AC^2 - 2 \cdot AE_4 \cdot AC \cdot \cos(\angle E_4AC).\]Substituting the known lengths and angles gives
 \[CE_4^2 = 222 - 222 \cdot \cos(120 + \alpha).\]Expanding this with the Cosine Addition Identity we get
 \[CE_4^2 = 222 - 222(\cos120 \cos \alpha - \sin120 \sin \alpha).\]Adding this to our equation for $CE_2^2$, we get
 \[CE_2^2 + CE_4^2 = 444 - 222(\cos120 \cos \alpha - \sin120 \sin \alpha) - 222(\cos120 \cos \alpha + \sin120 \sin \alpha).\]Simplifying we get
 \[CE_2^2 + CE_4^2 = 444 - 444(\cos120 \cos \alpha)\]

 We can find $\cos \alpha$ by using Law of Cosines on $\bigtriangleup BAD_1$. This gives
 \[11 = 222 - 222 \cos \alpha.\]Thus $\cos \alpha = \frac{211}{222}$.Substituting it in gives
 \[CE_2^2 + CE_4^2 = 444 - 444(\cos120 \cdot \frac{211}{222}).\]Thus
 \[CE_2^2 + CE_4^2 = 444 + 211 = 655.\]

 Therefore the desired sum is equal to 

 
 \[11+11+655 = \framebox{677}.\]

 



\textbf{Solution 6}


                    \begin{figure}[H]
                        \centering
                        \includegraphics[width=0.5\textwidth]{Solution/AIME_II_2012_13_Solution_Figure_1.png}
                    \end{figure}
                    

 We create a diagram. Each of the red quadrilaterals are actually 60 degree rotations about point A. We easily get that $\overline{C E_1} = \sqrt{11}$, and $\overline{C E_3} = \sqrt{11}$. If we set $\angle B A D_2 = \theta$, we can start angle chasing. In particular, we will like to find $\angle D E_4 C$, and $\angle D E_2 C$, since then we will be able to set up some Law Of Cosines. $\angle D E_4 C = \angle D E_4 A + \angle A E_4 C = 90 - \frac{\theta}{2} + 30 + \frac{\theta}{2} = 120^{\circ}$ That was convenient! We can do it with the other angle as well. $\angle D E_2 C = \angle D E_2 A - \angle C E_2 A = 90 - \frac{\theta}{2} - (30 - \frac{\theta}{2}) = 60^{\circ}$. That means we are able to set up Law of Cosines, on triangles $\triangle D E_4 C$ and $\triangle D E_2 C$, with some really convenient angles.Let $CE_2 = x$, and $CE_4 = y$.
 \[333 = 11 + x^2 - \sqrt{11} x\]
 \[333 = 11 + y^2 + \sqrt{11} y\]We subtract and get:
 \[0 = (x+y)(x-y-\sqrt{11})\] $x+y$ obviously can't be 0, so $x-y = \sqrt{11}$We add and get:
 \[666 = 22 + x^2 + y^2 + \sqrt{11} (y-x)\]. $y-x = -\sqrt{11}$. Thus, we can fill in and solve.
 \[666 = 22 + x^2 + y^2 - 11\]
 \[655 = x^2 + y^2\]Thus our answer is $C E_1^2 + C E_2^2 + C E_2^2 + C E_4^2 = 11 + 11 + C E_2^2 + C E_4^2 = 11 + 11 + x^2 + y^2 = 11 + 11 + 655 = \boxed{677}$.

 -Alexlikemath

 



\textbf{Solution 7}

Let $D_2, E_2, E_4$ be on the same side of $AB$ as $C.$ It is clear that $CE_2=CE_4=\sqrt{11}.$ Then, consider quadrilateral $CE_3BD_2.$ Since $\triangle BD_2E_3$ and $\triangle D_2BC$ are congruent, $D_2BE_3C$ is cyclic because $\angle BE_3D_2\cong\angle BCD_2.$ By Ptolemy's Theorem, we have $CE_3=\frac{111-CD_2^2}{\sqrt{11}}.$ Similarly, $CE_1=\frac{CD_1^2-111}{\sqrt{11}}.$

 Let $\angle ABD_2\cong\angle AD_2B=\theta.$ Then, $\cos\theta=\frac{\sqrt{11}}{2\sqrt{111}},$ so $\sin\theta=\frac{\sqrt{433}}{2\sqrt{111}}.$ Now, LoC on $\triangle BD_2C$ gives $CD_2^2=11+111-2\sqrt{11}\sqrt{111}\cos (\theta-60^\circ)$ and LoC on $\triangle BD_1C$ gives $CD_1^2=11+111-2\sqrt{11}\sqrt{111}\cos (\theta+60^\circ).$ By the cosine sum rule, we get\begin{align*}CE_3&=\frac{111-CD_2^2}{\sqrt{11}}\\&=\frac{2\sqrt{11}\sqrt{111}\cos (\theta-60^\circ)-11}{\sqrt{11}}\\&=2\sqrt{111}(\cos\theta\cos 60^\circ+\sin\theta\sin 60^\circ)-\sqrt{11}\\&=2\sqrt{111}\left(\frac{\sqrt{11}+\sqrt{1299}}{4\sqrt{111}}\right)-\sqrt{11}\\&=\frac{\sqrt{1299}-\sqrt{11}}{2}.\\\end{align*}
Similarly,\begin{align*}CE_1&=\frac{CD_1^2-111}{\sqrt{11}}\\&=\frac{11-2\sqrt{11}\sqrt{111}\cos (\theta+60^\circ)}{\sqrt{11}}\\&=\sqrt{11}-2\sqrt{111}(\cos\theta\cos 60^\circ-\sin\theta\sin 60^\circ)\\&=\sqrt{11}-2\sqrt{111}\left(\frac{\sqrt{11}-\sqrt{1299}}{4\sqrt{111}}\right)\\&=\frac{\sqrt{1299}+\sqrt{11}}{2}.\\\end{align*}
Therefore, our desired answer is\begin{align*}CE_1^2+CE_2^2+CE_3^2+CE_4^2&=CE_1^2+CE_3^2+11+11\\&=2\left(\frac{\sqrt{1299}^2+\sqrt{11}^2}{2^2}\right)+22\\&=\frac{1310}{2}+22\\&=\boxed{677}.\\\end{align*}


 -pieMax2713

 



\textbf{Solution 8}


                    \begin{figure}[H]
                        \centering
                        \includegraphics[width=0.5\textwidth]{Solution/AIME_II_2012_13_Solution_Figure_2.png}
                    \end{figure}
                    

 WLOG assume $CD_2>CD_1, CE_2>CE_1, CE_4>CE_3$.

 We first know that $CE_1=CE_3=\sqrt{11}$.

 The supplementary angle of $\angle E_1CE_2$ is half the value of $\angle E_2AE_1$. Also, $\angle E_2AE_1=\angle D_1AE_1 + \angle D_1AE_2 = 120^{\circ}$, so $\angle E_1CE_2=180^{\circ}-\frac{120^{\circ}}{2}=120^{\circ}$.

 We now use Law of Cosines on $\Delta CE_1E_2$ and $\angle E_1CE_2$:
 \[(E_1E_2)^2=(CE_1)^2+(CE_2)^2-2(CE_1)(CE_2)\cos(120^{\circ})\]
 \[333=(CE_2)^2+11-2(CE_2)(\sqrt{11})(-\frac{1}{2})\]
 \[333=(CE_2)^2-\sqrt{11}(CE_2)+11\]
 \[CE_2=\frac{\sqrt{1299}-\sqrt{11}}{2}\]

 
                    \begin{figure}[H]
                        \centering
                        \includegraphics[width=0.5\textwidth]{Solution/AIME_II_2012_13_Solution_Figure_3.png}
                    \end{figure}
                    

 Let $F$ be the midpoint of $E_3C$. Since $\Delta AE_3C$ is isosceles, $AF$ is the altitude of $E_4C$. Then
 \[\angle E_3E_4C=\frac{1}{2}\angle E_3AC=\angle E_3AF\]
 \[\sin(\angle E_3E_4C)=\sin(\angle E_3AF)=\frac{\frac{\sqrt{11}}{2}}{\sqrt{111}}=\frac{\sqrt{11}}{2\sqrt{111}}\]
 \[\cos(\angle E_3E_4C)=\sqrt{1-\sin(\angle E_3E_4C)^2}=\frac{\sqrt{433}}{2\sqrt{111}}\]

 Now the Law of Cosines on $\Delta CE_3E_4$ and $\angle E_3E_4C$:
 \[(CE_3)^2=(CE_4)^2+(E_3E_4)^2-2(CE_4)(E_3E_4)\cos(\angle E_3E_4C)\]
 \[11=(CE_4)^2+333-2(CE_4)(\sqrt{333})(\frac{\sqrt{433}}{2\sqrt{111}})\]
 \[11=(CE_4)^2-\sqrt{1299}(CE_4)+333\]
 \[CE_4=\frac{\sqrt{1299}+\sqrt{11}}{2}\]

 Finally,
 \[\sum_{k=1}^4(CE_k)^2\]
 \[=((CE_2)^2+(CE_4)^2)+(CE_1)^2+(CE_3)^2\]
 \[=(655)+(11)+(11)\]
 \[=\framebox{677}\]

 -Marchk26

 



\textbf{Solution 9 (Coordinate Geometry)}

Note that $D_1,D_2$ are the intersection of the circle with radius $\sqrt{11}$ centered at $B$ and the circle with radius $\sqrt{111}$ centered at $A$. Let these circle be $\omega, \Omega$. We establish a complex plane with $B$ being the origin. Here is out setup:\begin{align*}A&=\left(\frac{\sqrt{111}}{2},\frac{3\sqrt{37}}{2}\right)\\B&=(0,0)\\C&=(\sqrt{111},0)\\D_1&=(x_1,y_1)\\D_2&=(x_2,y_2)\\\omega &: x^2+y^2=11\\\Omega &: \left(x-\frac{\sqrt{111}}{2}\right)^2+\left(y-\frac{3\sqrt{37}}{2}\right)^2=111\end{align*}
The coordinate of $E_1$ is the result of a $60^\circ$ rotation of $A-D_1$ about $D_1$. We find\begin{align*}E_1&=x_1+y_1i+\left(\Bigl(\frac{\sqrt{111}}{2}-x_1\Bigr)+i\Bigl(\frac{3\sqrt{37}}{2}-y_1\Bigr)\right)\left(\frac{1}{2}+\frac{\sqrt{3}}{2}i\right)\\&=\Biggl(x_1+\frac{1}{2}\Bigl(\frac{\sqrt{111}}{2}-x_1\Bigr)-\frac{\sqrt{3}}{2}\Bigl(\frac{3\sqrt{37}}{2}-y_1\Bigr)\Biggr) \\&\qquad +i\Biggl(y_1+\frac{1}{2}\Bigl(\frac{3\sqrt{37}}{2}-y_1\Bigr)+\frac{\sqrt{3}}{2}\Bigl(\frac{\sqrt{111}}{2}-x_1\Bigr)\Biggr) \\&=\frac{1}{2}x_1+\frac{\sqrt{3}}{2}y_1-\frac{\sqrt{111}}{2}+i\left(\frac{1}{2}y_1-\frac{\sqrt{3}}{2}x_1+\frac{3\sqrt{37}}{2}\right)\end{align*}
Similarly, we obtain the coordinate of $E_2$ by a $-60^\circ$ rotation of $A-D_1$. We have
 \[E_2=\frac{1}{2}x_1-\frac{\sqrt{3}}{2}y_1+\sqrt{111}+i\left(\frac{1}{2}y_1+\frac{\sqrt{3}}{2}x_1\right)\]By the distance formula, we have\begin{align*}CE_1^2+CE_2^2&=\left(\frac{1}{2}x_1+\frac{\sqrt{3}}{2}y_1-\frac{3\sqrt{111}}{2}\right)^2+\left(\frac{1}{2}y_1-\frac{\sqrt{3}}{2}x_1+\frac{3\sqrt{37}}{2}\right)^2\\&\qquad+\left(\frac{1}{2}x_1-\frac{\sqrt{3}}{2}y_1\right)^2+\left(\frac{1}{2}y_1+\frac{\sqrt{3}}{2}x_1\right)^2\end{align*}
\begin{align*}CE_1^2+CE_2^2&=\frac{1}{2}x_1^2+\frac{3}{2}y_1^2-3\sqrt{111}\left(\frac{1}{2}x_1+\frac{\sqrt{3}}{2}y_1\right)+\frac{999}{4}+\frac{333}{4}+\frac{1}{2}y_1^2+\frac{3}{2}x_1^2\\&=2x_1^2+2y_1^2-3\sqrt{37}y_1-3\sqrt{111}x_1+333\end{align*}
By symmetry, we have\begin{align*}\sum_{k=1}^{4}(CE_k)^2&=\sum_{k=1}^{2}\left(2x_k^2+2y_k^2-3\sqrt{37}y_k-3\sqrt{111}x_k+333\right)\\&=710-3\sqrt{37}(y_1+y_2)-3\sqrt{111}(x_1+x_2)\tag{1}\end{align*}
Subtracting the equation of $\omega$ from the equation of $\Omega$, we see that 
 \[-\sqrt{111}x-3\sqrt{37}y+\frac{111}{4}+\frac{333}{4}=100\iff \sqrt{111}x+3\sqrt{37}y=11\]Obtaining the expression for $x^2,y^2$ and substituting them back to the equation of $\omega$, we have\begin{gather*}\frac{4}{3}x^2-\frac{22\sqrt{111}}{333}x+C_1=0\\4y^2-\frac{66\sqrt{37}}{111}y+C_2=0\end{gather*}Hence, by Vieta's formula,
 \[x_1+x_2=\frac{11\sqrt{111}}{222},\quad y_1+y_2=\frac{33\sqrt{37}}{222}\]By $(1)$, we have
 \[\sum_{k=1}^{4}(CE_k)^2=710-3\sqrt{37}\cdot \frac{33\sqrt{37}}{222}-3\sqrt{111}\cdot\frac{11\sqrt{111}}{222}=\boxed{677}\]

 


	\item (\textit{AIME II 2012}) In a group of nine people each person shakes hands with exactly two of the other people from the group. Let $N$ be the number of ways this handshaking can occur. Consider two handshaking arrangements different if and only if at least two people who shake hands under one arrangement do not shake hands under the other arrangement. Find the remainder when $N$ is divided by $1000$.


 
 \textbf{Answer: }016



\textbf{Solution 1}

Given that each person shakes hands with two people, we can view all of these through graph theory as 'rings'.  This will split it into four cases:  Three rings of three, one ring of three and one ring of six, one ring of four and one ring of five, and one ring of nine. (All other cases that sum to nine won't work, since they have at least one 'ring' of two or fewer points, which doesn't satisfy the handshaking conditions of the problem.)

 \textbf{Case 1:}   To create our groups of three, there are $\dfrac{\dbinom{9}{3}\dbinom{6}{3}\dbinom{3}{3}}{3!}$.  In general, the number of ways we can arrange people within the rings to count properly is $\dfrac{(n-1)!}{2}$, since there are $(n-1)!$ ways to arrange items in the circle, and then we don't want to want to consider reflections as separate entities.  Thus, each of the three cases has $\dfrac{(3-1)!}{2}=1$ arrangements.  Therefore, for this case, there are $\left(\dfrac{\dbinom{9}{3}\dbinom{6}{3}\dbinom{3}{3}}{3!}\right)(1)^3=280$

 \textbf{Case 2:}  For three and six, there are $\dbinom{9}{6}=84$ sets for the rings.  For organization within the ring, as before, there is only one way (i.e. $(3-1)!/2$)to arrange the ring of three.  For six, there is $\dfrac{(6-1)!}{2}=60$.  This means there are $(84)(1)(60)=5040$ arrangements.

 \textbf{Case 3:}  For four and five, there are $\dbinom{9}{5}=126$ sets for the rings.  Within the five, there are $\dfrac{4!}{2}=12$, and within the group of four there are $\dfrac{3!}{2}=3$ arrangements.  This means the total is $(126)(12)(3)=4536$.

 \textbf{Case 4:}  For the nine case, there is $\dbinom{9}{9}=1$ arrangement for the ring.  Within it, there are $\dfrac{8!}{2}=20160$ arrangements.

 Summing the cases, we have $280+5040+4536+20160=30016 \to \boxed{016}$.

 Note: the reason why we divided by $3!$ in Case 1 and not the other three cases is because for case one and case one only, the groups each have 3 people, implying a symmetric configuration. 

 For example, for 3 people A, B, C in group 1, they could very well be in group 2 or 3 too, because each group has 3 people. 

 However, for all the other three cases, the configurations are not symmetric and so dividing by $2!$ or like would be incorrect. It is not possible to configure A, B, C such that they can be both a group of 3 and a group of 6, since {A, B, C} only has 3 elements. 

 



\textbf{Solution 2}

Let $f_N$ be the number of ways a group of $N$ people can shake hands with exactly two of the other people from the group, where $N \ge 3$.

 We can easily find that $f_3=1$.

 Continuing on, we will label the people as $A$,$B$,$C$,... corresponding with $N$.

 $f_4$: There are $\dbinom{3}{2}=3$ possible ways for person $A$ to shake hands with two others (WLOG assume $A$ shakes hands with $B$ and $C$), and there is only one possible outcome thereon ($B$ and $C$ both shakes hands with $D$).

 Therefore, we can conclude that $f_4=3$

 $f_5$: There are $\dbinom{4}{2}=6$ possible ways for person $A$ to shake hands with two others (WLOG assume they are $B$ and $C$). Then, $B$ and $C$ must also shake hands with the remaining two people ($2$ ways), and those last $2$ people must shake hands with each other ($1$ way).Therefore, $f_5=\dbinom{4}{2}\ast(2)\ast(1)=12$

 $f_6$: There are $\dbinom{5}{2}=10$ possible ways for person $A$ to shake hands with two others (WLOG assume they are $B$ and $C$). However, $B$ and $C$ shake hands with each other, then there are $f_3$ ways for the rest of the people to shake hands. If $B$ and $C$ shake hands with two others out of the three remaining people (3$\ast$2 ways), those $2$ people must shake hands with the last person ($1$ way).Therefore, $f_6=\dbinom{5}{2}\ast(f_3)+3\ast(2)\ast(1)=70$

 Now we have enough information to find $f_9$

 $f_9$: There are $\dbinom{8}{2}=28$ possible ways for person $A$ to shake hands with two others (WLOG assume they are $B$ and $C$).

 \textbf{Case 1:} $B$ and $C$ shake hands with each other

 There are $f_6$ ways for $D$,$E$,$F$,$G$,$H$, and $I$ to shake hands

 \textbf{Case 2:} $B$ and $C$ shake hands with one of the others out of the $6$ remaining people ($6$ ways)

 There are $f_5$ ways for $E$,$F$,$G$,$H$, and $I$ to shake hands

 \textbf{Case 3:} If $B$ and $C$ shake hands with two different people (there are 6$\ast$5 ways to choose the people but WLOG assume they are $D$ and $E$). Calculations of Case 3 are in its subcases.

 \textbf{Subcase 3.1:} $D$ and $E$ shake hands with each otherThere are $f_4$ ways for $F$,$G$,$H$, and $I$ to shake hands

 \textbf{Subcase 3.2:} $D$ and $E$ each shake hands with one other person (there are $4$ ways to choose the person, WLOG let that be $F$)There are $f_3$ ways for $G$,$H$, and $I$ to shake hands

 
\textbf{Subcase 3.3:} $D$ and $E$ each shake hands with two different people (there are 4$\ast$3 ways, WLOG let that be $F$ and $G$)There are $2\ast1$ ways for $F$ and $G$ to then shake hands with $H$ and $I$, and $H$ and $I$ will shake hands with each other.

 We have 

 
 \begin{align*} 28(f_6+6\ast(f_5)+6\ast(5)\ast(f_4)+4\ast(f_3)+4\ast(3)\ast(2)\ast(1) &= 28(70+6\ast(12)+6\ast(5)\ast(3+4\ast(1)+4\ast(3)\ast(2)\ast(1) \\ &= 30016 \end{align*}


 which gives us the answer of $\boxed{016}$

 


	\item (\textit{AIME II 2012}) Triangle $ABC$ is inscribed in circle $\omega$ with $AB=5$, $BC=7$, and $AC=3$. The bisector of angle $A$ meets side $\overline{BC}$ at $D$ and circle $\omega$ at a second point $E$. Let $\gamma$ be the circle with diameter $\overline{DE}$. Circles $\omega$ and $\gamma$ meet at $E$ and a second point $F$. Then $AF^2 = \frac mn$, where $m$ and $n$ are relatively prime positive integers. Find $m+n$.


 
 \textbf{Answer: }919



\textbf{Solution 1}

Let $\tau$ be $\sqrt{bc}$-inversion. $\tau$ sends $B \leftrightarrow C$ and $D\leftrightarrow E$. Since $DE$ is the diameter of $\gamma$, this implies $\gamma$ stays fixed. If $M$ is the midpoint of $\overline{BC}$, 
 \[\tau(M) = \tau(\overline{BC} \cap \gamma) = \tau(\overline{BC})\cap \tau(\gamma) = \omega \cap \gamma = F\] We conclude $AF\cdot AM = bc = 15$. Now, switching to barycentrics, $\overrightarrow{AM} = (-1, \frac{1}{2}, \frac{1}{2})$, so we find 
 \[AM^2 = -a^2yz-b^2xz-c^2xy = -\frac{49}{4} +\frac{9}{2}+\frac{25}{2} = \frac{19}{4} \implies AF^2 = \frac{225}{\frac{19}{4}}=\frac{900}{19} \implies \boxed{919}\]

 -sillybone

 \textbf{Quick Solution using Olympiad Terms} Take a force-overlaid inversion about $A$ and note $D$ and $E$ map to each other. As $DE$ was originally the diameter of $\gamma$, $DE$ is still the diameter of $\gamma$. Thus $\gamma$ is preserved. Note that the midpoint $M$ of $BC$ lies on $\gamma$, and $BC$ and $\omega$ are swapped. Thus points $F$ and $M$ map to each other, and are isogonal. It follows that $AF$ is a symmedian of $\triangle{ABC}$, or that $ABFC$ is harmonic. Then $(AB)(FC)=(BF)(CA)$, and thus we can let $BF=5x, CF=3x$ for some $x$. By the LoC, it is easy to see $\angle{BAC}=120^\circ$ so $(5x)^2+(3x)^2-2\cos{60^\circ}(5x)(3x)=49$. Solving gives $x^2=\frac{49}{19}$, from which by Ptolemy's we see $AF=\frac{30}{\sqrt{19}}$. We conclude the answer is $900+19=\boxed{919}$.

 \textbf{- Emathmaster}

 Side Note: You might be wondering what the motivation for this solution is. Most of the people who've done EGMO Chapter 8 should recognize this as problem 8.32 (2009 Russian Olympiad) with the computational finish afterwards.Now if you haven't done this, but still know what inversion is, here's the motivation. We'd see that it's kinda hard to angle chase, and if we could, it would still be a bit hard to apply (you could use trig, but it won't be so clean most likely). If you give up after realizing that angle chasing won't work, you'd likely go in a similar approach to Solution 2 (below) or maybe be a bit more insightful and go with the elementary solution above.

 Finally, we notice there's circles! Classic setup for inversion! Since we're involving an angle-bisector, the first thing that comes to mind is a force overlaid inversion described in Lemma 8.16 of EGMO (where we invert with radius $\sqrt{AB \cdot AC}$ and center $A$, then reflect over the $A$-angle bisector, which fixes $B, C$). We try applying this to the problem, and it's fruitful - we end up with this solution.-MSC

 



\textbf{Solution 2}

Use the angle bisector theorem to find $CD=\frac{21}{8}$, $BD=\frac{35}{8}$, and use Stewart's Theorem to find $AD=\frac{15}{8}$. Use Power of Point $D$ to find $DE=\frac{49}{8}$, and so $AE=8$. Use law of cosines to find $\angle CAD = \frac{\pi} {3}$, hence $\angle BAD = \frac{\pi}{3}$ as well, and $\triangle BCE$ is equilateral, so $BC=CE=BE=7$.
                    \begin{figure}[H]
                        \centering
                        \includegraphics[width=0.5\textwidth]{Solution/AIME_II_2012_15_Solution_Figure_1.png}
                    \end{figure}
                    In triangle $AEF$, let $X$ be the foot of the altitude from $A$; then $EF=EX+XF$, where we use signed lengths. Writing $EX=AE \cdot \cos \angle AEF$ and $XF=AF \cdot \cos \angle AFE$, we get
 \begin{align}\tag{1}     EF = AE \cdot \cos \angle AEF + AF \cdot \cos \angle AFE. \end{align}Note $\angle AFE = \angle ACE$, and the Law of Cosines in  $\triangle ACE$ gives  $\cos \angle ACE = -\frac 17$. Also, $\angle AEF = \angle DEF$, and $\angle DFE = \frac{\pi}{2}$ ($DE$ is a diameter), so $\cos \angle AEF = \frac{EF}{DE} =  \frac{8}{49}\cdot EF$.

 Plugging in all our values into equation $(1)$, we get:
 \[EF = \frac{64}{49} EF -\frac{1}{7} AF \quad \Longrightarrow \quad EF = \frac{7}{15} AF.\]The Law of Cosines in  $\triangle AEF$, with $EF=\frac 7{15}AF$ and $\cos\angle AFE = -\frac 17$ gives
 \[8^2 = AF^2 + \frac{49}{225}  AF^2 + \frac 2{15} AF^2 = \frac{225+49+30}{225}\cdot AF^2\]Thus $AF^2 = \frac{900}{19}$. The answer is $\boxed{919}$.~Shen Kislay Kai

 



\textbf{Solution 3}

Let $a = BC$, $b = CA$, $c = AB$ for convenience. Let $M$ be the midpoint of segment $BC$. We claim that $\angle MAD=\angle DAF$. 

 $\textit{Proof}$. Since $AE$ is the angle bisector, it follows that $EB = EC$ and consequently $EM\perp BC$. Therefore, $M\in \gamma$. Now let $X = FD\cap \omega$. Since $\angle EFX=90^\circ$, $EX$ is a diameter, so $X$ lies on the perpendicular bisector of $BC$; hence $E$, $M$, $X$ are collinear. From $\angle DAG = \angle DMX = 90^\circ$, quadrilateral $ADMX$ is cyclic. Therefore, $\angle MAD = \angle MXD$. But $\angle MXD$ and $\angle EAF$ are both subtended by arc $EF$ in $\omega$, so they are equal. Thus $\angle MAD=\angle DAF$, as claimed. 
                    \begin{figure}[H]
                        \centering
                        \includegraphics[width=0.5\textwidth]{Solution/AIME_II_2012_15_Solution_Figure_2.png}
                    \end{figure}
                    As a result, $\angle CAM = \angle FAB$. Combined with $\angle BFA=\angle MCA$, we get $\triangle ABF\sim\triangle AMC$ and therefore
 \[\frac c{AM}=\frac {AF}b\qquad \Longrightarrow \qquad AF^2=\frac{b^2c^2}{AM^2} = \frac{15^2}{AM^2}\]By Stewart's Theorem on $\triangle ABC$ (with cevian $AM$), we get 
 \[AM^2 = \frac 12 (b^2+c^2)-\frac 14 a^2 = \frac{19}{4},\] so $AF^2 = \frac{900}{19}$, so the answer is $900+19=\boxed{919}$. 

 \textbf{-Solution by thecmd999}

 



\textbf{Solution 4}

Use the angle bisector theorem to find $CD=\frac{21}{8}$, $BD=\frac{35}{8}$, and use Stewart's Theorem to find $AD=\frac{15}{8}$. Use Power of Point $D$ to find $DE=\frac{49}{8}$, and so $AE=8$. Then use the Extended Law of Sine to find that the length of the circumradius of $\triangle ABC$ is $\frac{7\sqrt{3}}{3}$. 
                    \begin{figure}[H]
                        \centering
                        \includegraphics[width=0.5\textwidth]{Solution/AIME_II_2012_15_Solution_Figure_3.png}
                    \end{figure}
                    Since $DE$ is the diameter of circle $\gamma$, $\angle DFE$ is $90^\circ$. Extending $DF$ to intersect circle $\omega$ at $X$, we find that $XE$ is the diameter of $\omega$ (since $\angle DFE$ is $90^\circ$). Therefore, $XE=\frac{14\sqrt{3}}{3}$.

 Let $EF=x$, $XD=u$, and $DF=v$. Then $XE^2-XF^2=EF^2=DE^2-DF^2$, so we get
 \[(u+v)^2-v^2=\frac{196}{3}-\frac{2401}{64}\]which simplifies to
 \[u^2+2uv = \frac{5341}{192}.\]By Power of Point $D$, $uv=BD \cdot DC=735/64$. Combining with above, we get
 \[XD^2=u^2=\frac{931}{192}.\] Note that $\triangle XDE\sim \triangle ADF$ and the ratio of similarity is $\rho = AD : XD = \frac{15}{8}:u$. Then $AF=\rho\cdot XE = \frac{15}{8u}\cdot R$ and 
 \[AF^2 = \frac{225}{64}\cdot \frac{R^2}{u^2} = \frac{900}{19}.\]The answer is $900+19=\boxed{919}$. 

 \textbf{-Solution by TheBoomBox77}

 



\textbf{Solution 5}

Use Law of Cosines in $\triangle ABC$ to get $\angle BAC=120^\circ$. Because $AE$ bisects $\angle A$, $E$ is the midpoint of major arc $BC$ so $BE=CE,$ and $\angle BEC=60^\circ.$ Thus $\triangle BEC$ is equilateral. Notice now that $\angle BFC=\angle BFE= 60^\circ.$ But $\angle DFE=90^\circ$ so $FD$ bisects $\angle BFC.$ Thus, 
 \[\frac{BF}{CF}=\frac{BD}{CD}=\frac{BA}{CA}=\frac{5}{3}.\]Let $BF=5k, CF=3k.$ Use Law of Cosines on $\triangle BFC$ to get
 \[25k^2+9k^2-15k^2 = 49 \qquad \Longrightarrow\qquad k=\frac 7{\sqrt{19}}\] Use Ptolemy's Theorem on $BFCA$, to get 
 \[15k+15k=7\cdot AF, \qquad \Longrightarrow\qquad AF= \frac{30}{\sqrt{19}},\] so $AF^2=\frac{900}{19}$ and the answer is $900+19=919$ 

 



\textbf{Solution 6}


                    \begin{figure}[H]
                        \centering
                        \includegraphics[width=0.5\textwidth]{Solution/AIME_II_2012_15_Solution_Figure_4.png}
                    \end{figure}
                    Denote $AB = c, BC = a, AC = b, \angle A = 2 \alpha.$Let M be midpoint BC. Let $\theta$ be the circle centered at $A$ with radius $\sqrt{AB \cdot AC} =\sqrt{bc}.$

 We calculate the length of some segments.The median $AM = \sqrt{\frac {b^2}{2} + \frac {c^2}{2} -  \frac {a^2}{4}}.$  The bisector $AD = \frac {2 b c \cos \alpha}{b+c}.$ One can use Stewart's Theorem in both cases.

 $AD$ is bisector of $\angle A \implies BD = \frac {a c}{b + c}, CD = \frac {a b}{b + c} \implies$
 \[BD \cdot CD = \frac {a^2 bc }{(b+c)^2}.\]We use Power of Point $D$ and get $AD \cdot DE = BD \cdot CD.$
 \[AE = AD + DE = AD + \frac {BD \cdot CD}{AD},\]
 \[AE =\frac {2 b c \cos \alpha}{b+c} + \frac {a^2 bc \cdot (b+c) }{(b+c)^2 \cdot 2 b c \cos \alpha} =\]
 \[=  \frac {b c \cos^2 \alpha + a^2}{2(b+c)\cos \alpha} =\frac {4bc \cos^2 \alpha  + b^2 +c^2 -2 b c \cos 2\alpha}{2(b+c) \cos \alpha} = \frac {b+c}{2} \implies AD \cdot AE = 2 bc \cos \alpha.\]We consider the inversion with respect $\theta.$

 $B$ swap $B' \implies AB' = AC, B'  \in AB \implies B'$ is symmetric to $C$ with respect to $AE.$

 $C$ swap $C' \implies AC' = AB, C'$ lies on line $AC \implies C'$ is symmetric to $B$ with respect to $AE.$

 $BC^2 = AB^2 + AC^2 + AB \cdot BC \implies \alpha = 60^\circ \implies AD \cdot AE = bc \implies D$ swap $E.$

 Points $D$ and $E$ lies on $\Gamma \implies \Gamma$ swap $\Gamma.$

 $DE$ is diameter $\Gamma, \angle DME = 90^\circ \implies  M \in \Gamma.$ Therefore $M$ is crosspoint of $BC$ and $\Gamma.$

 Let $\Omega$ be circumcircle $AB'C'.  \Omega$ is image of line $BC.$Point $M$ maps into $M' \implies M' = \Gamma \cap \Omega.$

 Points $A, B',$ and $C'$ are symmetric to $A, C,$ and $B,$ respectively. 

 Point $M'$ lies on $\Gamma$ which is symmetric with respect to $AE$ and on $\Omega$ which is symmetric to $\omega$ with respect to $AE  \implies$

 <i>\textbf{$M'$ is symmetric $F$ with respect to $AE \implies AM' = AF.$}</i>

 We use Power of Point $A$ and get 
 \[AF = AM' = \frac {AD \cdot AE}{AM} = \frac {4b c}{\sqrt{2 b^2 + 2 c^2 – a^2}} = \frac {4 \cdot 3 \cdot 5}{\sqrt{ 50 + 18 – 49}} = \frac {30}{\sqrt{19}} \implies \boxed{\textbf{919}}.\]

 \textbf{vladimir.shelomovskii@gmail.com, vvsss}

 



\textbf{Solution 7}

To do this, we first define the intersection of $EF$ and $BC$ to be $K$. 

 \textbf{Lemma 1:} $(K, C, D, B)$ are harmonic. First of all, define the midpoint of $BC$ to be $M$. Then, we have that angle FMD is $90$ degrees, and as a result, $M$ lies on this circle. By \href{https://artofproblemsolving.com/wiki/index.php?title=Power\_of\_a\_Point}{Power of a Point}, $(KD)(KM) = (KE)(KF) = (KB)(KC)$. As a result, $(K, C, D, B)$ are harmonic from another famous harmonic lemma.

 As a result, since $\angle EFD = 90^{\circ}$, by another Harmonic Lemma, $FD$ is the angle bisector of $BFC$. Since $\frac{BD}{CD} = \frac{5}{3}$ by \href{https://artofproblemsolving.com/wiki/index.php?title=Angle\_bisector\_theorem}{angle bisector theorem}, $\frac{BF}{CF} = \frac{5}{3}$. Since $\angle BAC$ is $120^{\circ}$ by \href{https://artofproblemsolving.com/wiki/index.php?title=Law\_of\_Cosines}{Law of Cosines} (LOC), we can use LOC to finish off. Call $BF = 5a$, and $CF = 3a$, $(5a)^2+(3a)^2-15a^2 = 49$, so $a = \frac{7}{\sqrt{19}}$. We do \href{https://artofproblemsolving.com/wiki/index.php?title=Ptolemy\%27s\_Theorem}{Ptolemy's Theorem} on $ABFC$. Our answer is: 
 \[\frac{(3)(35)+(5)(21)}{7\sqrt{19}} = \frac{30}{\sqrt{19}}.\]As a result, the final answer is $\boxed{919}$.

 -sepehr2010

 Minor edits ~\href{https://artofproblemsolving.com/wiki/index.php?title=Zhenghua}{Zhenghua}

 



\textbf{Solution 8 (Coordinate Geometry)}

By the Angle Bisector Theorem, $CD=\frac{21}{8}$. Let $H$ be the foot of the perpendicular from $A$ to $\overline{BC}$. By the Pythagorean Theorem, $CH=\frac{33}{14}$, $AH=\frac{15\sqrt{3}}{14}$. By the Law of Cosines and the Law of Sines, 
 \[\cos A=-\frac{1}{2},\quad \sin A=\frac{\sqrt{3}}{2},\quad R=\frac{7\sqrt{3}}{3}.\]  Let $M$ be the midpoint of $BC$. Let $O_1,O_2$ be the centers of circles $\omega,\gamma$, respectively. We know that the point $E$ lies on $MO_1$. Since $\angle MO_1C=60^\circ$, we have $MO_1=\frac{7\sqrt{3}}{6}$. We establish a coordinate system with $O_1$ as the origin and $\overline{BC}$ parallel to the $x$-axis. After some minor calculations, here is our setup:\begin{align*}A&=\left(-\frac{8}{7},\frac{47\sqrt{3}}{21}\right)\\D&=\left(-\frac{7}{8},\frac{7\sqrt{3}}{6}\right)\\E&=\left(0,-\frac{7\sqrt{3}}{3}\right)\\F&=(x_1,y_1)\\M&=\left(0,\frac{7\sqrt{3}}{6}\right)\\O_1&=(0,0)\\O_2&=\left(-\frac{7}{16},-\frac{7\sqrt{3}}{12}\right)\\\omega&: x^2+y^2=\frac{49}{3}\\\gamma&: \left(x+\frac{7}{16}\right)^2+\left(y+\frac{7\sqrt{3}}{12}\right)^2=\frac{2401}{256}\end{align*}
Subtracting the equation of $\omega$ from that of $\gamma$, we obtain\begin{gather*}\frac{7}{8}x+\frac{49}{256}+\frac{7\sqrt{3}}{6}y+\frac{49}{48}=-\frac{5341}{768}\quad \iff\\x+\frac{4\sqrt{3}}{3}y=-\frac{28}{3} \tag{1}\end{gather*}By $(1)$, we have\begin{gather*}\left(-\frac{28}{3}-\frac{4\sqrt{3}}{3}y\right)^2+y^2=\frac{49}{3}\\\frac{19}{3}y^2 + \frac{224\sqrt{3}}{9}y + C = 0\end{gather*}By Vieta's formulas, We have $-\frac{7\sqrt{3}}{3}+y_1 = -\frac{224\sqrt{3}}{57}\implies y_1=-\frac{91\sqrt{3}}{57}$. Solving for $x_1$ by $(1)$, we have $x_1=-\frac{56}{19}$.

 Thus, the coordinate of $F$ evaluates to $F=\left(-\frac{56}{19}, -\frac{91\sqrt{3}}{57}\right)$. By the distance formula,\begin{align*}AF^2&=\left(\frac{56}{19}-\frac{8}{7}\right)^2+\left(\frac{47\sqrt{3}}{21}+\frac{91\sqrt{3}}{57}\right)^2\\&=\left(\frac{240}{19\cdot 7}\right)^2+\left(\frac{510\sqrt{3}}{19\cdot 7}\right)^2\\&=\frac{2^8\cdot 3^2\cdot 5^2}{19^2\cdot 7^2}+\frac{2^2\cdot 5^2\cdot 3^3\cdot 17^2}{19^2\cdot 7^2}\\&=\frac{2^2\cdot 3^2\cdot 5^2}{19^2\cdot 7^2}(2^6+3\cdot 17^2)\\&=\frac{2^2\cdot 3^2\cdot 5^2}{19^2\cdot 7^2}\cdot 7\cdot 7\cdot 19\\&=\frac{900}{19}\end{align*}
The requested sum is $900+19=\boxed{919}$.

 


\end{enumerate}